\BourbakiExerciseGroup

\begin{enumerate}
\def\labelenumi{\arabic{enumi})}
\tightlist
\item
  设 A, B 是拓扑空间 X 的两个子集，使得 \(A \supset B\)。证明：
\end{enumerate}

\begin{enumerate}
\def\labelenumi{\alph{enumi})}
\item
  B 相对于 X 的内部含于 B 相对于 X 的子空间 A 的内部。给出一个这两个集合不同的例子。
\item
  B 相对于 A 的边界含于 A 与 B 相对于 X 的边界之交。给出一个这两个集合不同的例子。
\end{enumerate}

\begin{enumerate}
\def\labelenumi{\arabic{enumi})}
\setcounter{enumi}{1}
\tightlist
\item
  设 A 与 B 是拓扑空间 X 的任意两个子集。
\end{enumerate}

\begin{enumerate}
\def\labelenumi{\alph{enumi})}
\item
  证明 A 与 B 相对于 X 的内部之交含于 \(B \cap A\) 相对于 A 的内部。给出一个这两个集合不同的例子。
\item
  证明 A 与 B 相对于 X 的闭包之交包含 B ∩ A 相对于 A 的闭包。给出一个这两个集合不同的例子。
\end{enumerate}

\begin{enumerate}
\def\labelenumi{\arabic{enumi})}
\setcounter{enumi}{2}
\item
  子空间 \(\mathbf{A}\)（属于拓扑空间 \(\mathbf{X}\)）是离散的，当且仅当 \(\mathbf{A}\) 的每一点都是孤立的。
\item
  设 Y, Z 是拓扑空间 X 的两个子空间，使得 \(X = Y \cup Z\)，且使得 \(\overline{A} \cap B = \overline{B} \cap A = \varnothing\)，其中 \(A = Y \cap CZ\)，\(B = Z \cap CY\)。
\end{enumerate}

\begin{enumerate}
\def\labelenumi{\alph{enumi})}
\item
  证明：若 M 是 X 的任一子集，则 M 在 X 中的闭包是 \(M \cap Y\) 在 Y 中的闭包与 \(M \cap Z\) 在 Z 中的闭包之并。由此推出，若 \(M \cap Y\) 在 Y 中是闭的（相应地，开的）且 \(M \cap Z\) 在 Z 中是闭的（相应地，开的），则 M 在 X 中是闭的（相应地，开的）。
\item
  设 \(f\) 是 \(\mathbf{X}\) 到拓扑空间 \(\mathbf{X}'\) 中的一个映射。证明：若 \(f|_{\mathbf{Y}}\) 与 \(f|_{\mathbf{Z}}\) 连续，则 \(f\) 也连续。
\end{enumerate}

\begin{enumerate}
\def\labelenumi{\arabic{enumi})}
\setcounter{enumi}{4}
\item
  设 Y, Z 是拓扑空间 X 的两个子空间，使得 \(X = Y \cup Z\)。证明：若 \(Y \cap Z\) 的子集 M 在 Y 与 Z 中都是开的（相应地，闭的），则 M 在 X 中是开的（相应地，闭的）。
\item
  设 A 是拓扑空间 X 的一个局部闭子空间。证明 X 的这样的开子集 U 所成的集——使得 \(A \subset U\) 且 A 在 U 中是闭的——有一个最大元素，即 A 相对于 \(\overline{A}\) 的边界在 X 中的余集。
\item
  用例子证明：两个局部闭子集的并与一个局部闭子集的余集未必是局部闭的。
\item
  设 A 是有序集 X 的一个子集。证明由 X 的右拓扑（相应地，左拓扑）（§ 1，习题 2）在 A 上诱导的拓扑，是 A 关于由 X 的序所诱导的 A 的序的右拓扑（相应地，左拓扑）。
\item
  设 \(\mathbf{A}\) 是线性序集 \(\mathbf{X}\) 的一个子集。
\end{enumerate}

\begin{enumerate}
\def\labelenumi{\alph{enumi})}
\item
  证明由拓扑 \(\mathcal{G}_{-}(\mathbf{X})\) {[}相应地 \(\mathcal{G}_{+}(\mathbf{X}),\mathcal{G}_{0}(\mathbf{X})]\)（§ 2，习题 5）在 A 上诱导的拓扑，比拓扑 \(\mathcal{G}_{-}(\mathbf{A})\) {[}相应地 \(\mathcal{G}_{+}(\mathbf{A}),\mathcal{G}_{0}(\mathbf{A})]\) 细，这里 A 的序是由 X 的序所诱导的序。若 A 是 X 的一个区间，则 \(\mathcal{G}_{-}(\mathbf{A})\) {[}相应地 \(\mathcal{G}_{+}(\mathbf{A}),\mathcal{G}_{0}(\mathbf{A})]\) 是由 \(\mathcal{G}_{-}(\mathbf{X})\) {[}相应地 \(\mathcal{G}_{+}(\mathbf{X}),\mathcal{G}_{0}(\mathbf{X})]\) 在 A 上诱导的拓扑。
\item
  取 X 为积 \(Q \times Z\)，按字典序排序（《集合论》，第三章，§ 2，第 6 目），并取 A 为子集 \(Q \times \{0\}\)。证明拓扑 \(\mathcal{G}_{-}(A)\) {[}相应地，\(\mathcal{G}_{+}(A)\)，\(\mathcal{G}_{0}(A)\){]} 严格地比由 \(\mathcal{G}_{-}(X)\) {[}相应地，\(\mathcal{G}_{+}(X)\)，\(\mathcal{G}_{0}(X)\){]} 在 A 上诱导的拓扑粗。
\end{enumerate}

\begin{enumerate}
\def\labelenumi{\arabic{enumi})}
\setcounter{enumi}{9}
\item
  设 A 是非空拓扑空间 X 的一个非空子空间。证明在 \(\mathfrak{P}_{0}(\mathbf{A})\) 上由拓扑 \(\mathcal{G}_{\Omega}\)（相应地，\(\mathcal{G}_{\Phi}\)）——\(\mathfrak{P}_{0}(\mathbf{X})\) 上的（§ 2，习题 7）——诱导的拓扑，就是拓扑 \(\mathcal{G}_{\Omega}\)（相应地，\(\mathcal{G}_{\Phi}\)）——\(\mathfrak{P}_{0}(\mathbf{A})\) 上的。
\item
  设 X 是拓扑空间，\(\mathcal{G}\) 是 X 的拓扑，A 是 X 的稠密子空间。证明 X 上比 \(\mathcal{G}\) 细、使 A 在 X 中稠密、且在 A 上诱导出与 \(\mathcal{G}\) 相同拓扑的诸拓扑所成的集至少有一个极大元素；这样的极大元素称为 A-极大拓扑。证明 X 上的拓扑 \(\mathcal{G}_0\) 是 A-极大的，当且仅当 X 的每个这样的子集 M——A∩M 在 M 中稠密（关于 \(\mathcal{G}_0\)）且在 A 中是开的——在拓扑 \(\mathcal{G}_0\) 中都是开的；且子空间 \(\mathbb{C}\) A 这时在由 \(\mathcal{G}_0\) 诱导的拓扑中是离散的，且在 X 中是闭的。再证明：若 \(\mathcal{G}_0\) 是 A-极大的，且由 \(\mathcal{G}_0\) 在 A 上诱导的拓扑是拟极大的（§ 2，习题 6），则 \(\mathcal{G}_0\) 是拟极大的。
\end{enumerate}

* 12) 设 X 是实直线 R 的子空间{[}0, 1{]}，S 是 X 上的等价关系，其等价类为 \{0, 1\} 以及诸集合 \{x\}（其中 0 \textless{} x \textless{} 1）。证明 X 关于 S 没有连续截面。*

\begin{enumerate}
\def\labelenumi{\arabic{enumi})}
\setcounter{enumi}{12}
\item
  设 X 是拓扑空间，R 与 S 是 X 上的两个等价关系，使得 R 蕴含 S。证明：若存在 X 关于 R 的连续截面，且存在 X/R 关于 S/R 的连续截面，则存在 X 关于 S 的连续截面。
\item
  设 X 是拓扑空间，它是两个子空间 \(X_{1}\) 与 \(X_{2}\) 的和，且设 X/R 是把 \(X_{1}\) 与 \(X_{2}\) 沿子空间 \(A_{1}\)（属于 \(X_{1}\)）与子空间 \(A_{2}\)（属于 \(X_{2}\)）借 \(A_{1}\) 到 \(A_{2}\) 上的一个同胚 h 粘合起来所得的空间。证明 \(X_{1}\) 与 \(X_{2}\) 在 X/R 中的典范像分别同胚于 \(X_{1}\) 与 \(X_{2}\)。
\end{enumerate}

* 15) 考虑 \(\mathbf{R}: \mathbf{X}_1 = ] - \mathbf{1}, \mathbf{1}[, \mathbf{X}_2 = ] - \mathbf{2}, -\mathbf{1}[, \mathbf{X}_3 = ]\mathbf{1}, \mathbf{2}[.\) 的以下三个子空间。设 \(\mathbf{X}\) 是它们的并；于是 \(\mathbf{X}\) 是这三个子空间的和。设 \(A_{2}\)（相应地，\(A_{3}\)）是含于 \(X_{2}\)（相应地，\(X_{3}\)）的无理数所成的集；设 \(B_{1}\)（相应地，\(B_{2}\)）是含于 {]}-1, 0{[}（相应地，\(X_{2}\)）的有理数所成的集；再设 \(C_{1}\)（相应地，\(C_{3}\)）是含于 {]}0, 1{[}（相应地，\(X_{3}\)）的有理数所成的集。设 X/R 是按下述方式把诸 \(X_{i}\) 粘合起来所得的空间：(i) 沿 \(A_{2}\) 与 \(A_{3}\) 借同胚 \(x \to x + 3\)；(ii) 沿 \(B_{1}\) 与 \(B_{2}\) 借同胚 \(x \to x - 1\)；(iii) 沿 \(C_{1}\) 与 \(C_{3}\) 借同胚 \(x \to x + 1\)。证明 \(X_{1}\) 在 X/R 中的典范像不同胚于 \(X_{1}\)。若 \(Y = X_{1} \cup B_{2} \cup C_{3}\) 是 \(X_{1}\) 关于 R 的饱和集，证明商空间 \(Y/R_{Y}\) 不同胚于 Y 在 X/R 中的典范像。*

\begin{enumerate}
\def\labelenumi{\arabic{enumi})}
\setcounter{enumi}{15}
\tightlist
\item
  设 X 是非空拓扑空间，R 是 X 上的一个等价关系。证明：若把 X/R 视为 \(\mathfrak{P}_{0}(X)\) 的一个子集，则 X/R 上的商拓扑比由拓扑 \(C_{\Omega}\) 与 \(C_{\Phi}\)——\(\mathfrak{P}_{0}(X)\) 上的（§ 2，习题 7）——在 X/R 上诱导的拓扑粗。
\end{enumerate}

